\documentclass[11pt]{article}
\usepackage[margin=1in]{geometry}
\usepackage{amsmath,amssymb,amsthm}
\usepackage[T1]{fontenc}
\usepackage{lmodern}
\usepackage{xcolor}
\usepackage{listings}
\usepackage{booktabs}
\usepackage[hidelinks]{hyperref}
\usepackage{url}

\newtheorem{theorem}{Theorem}
\newtheorem{lemma}[theorem]{Lemma}
\newtheorem{corollary}[theorem]{Corollary}
\newtheorem{proposition}[theorem]{Proposition}
\theoremstyle{remark}
\newtheorem{remark}[theorem]{Remark}

\newcommand{\leg}[2]{\left(\frac{#1}{#2}\right)}
\newcommand{\lean}[1]{\texttt{#1}}
% version-2 commits: primitive-root-families=b4ae30a94f5bc5bcba7e0143c239565af261038a
\newcommand{\repo}{https://github.com/jpbald93/primitive-root-families}
\newcommand{\commit}{e30457a}
\newcommand{\src}[2]{\href{\repo/blob/\commit/PrimitiveRootFamilies/#1\##2}{\texttt{#1}}}

\lstdefinelanguage{lean}{
  morekeywords={theorem,lemma,def,example,by,fun,have,show,exact,intro,namespace,end,import,private,local,instance,alias},
  sensitive=true, morecomment=[l]{--}, morecomment=[s]{/-}{-/}, morestring=[b]"}
\lstset{language=lean, basicstyle=\ttfamily\small, keywordstyle=\bfseries,
  columns=fullflexible, keepspaces=true, breaklines=true, frame=single, framerule=0.3pt,
  literate={ℕ}{{$\mathbb{N}$}}1 {ℤ}{{$\mathbb{Z}$}}1 {≠}{{$\neq$}}1 {↔}{{$\leftrightarrow$}}2 {∃}{{$\exists$}}1
           {∧}{{$\wedge$}}1 {→}{{$\to$}}1 {∣}{{$\mid$}}1 {≤}{{$\le$}}1 {⟨}{{$\langle$}}1 {⟩}{{$\rangle$}}1}

\title{Primitive roots modulo primes $2^k\ell+1$, formalised in Lean~4:\\
Chebyshev's base-10 criterion and the least primitive root of a safe prime}
\author{Josh Bald\\ \small Independent researcher, Ontario, Canada\\ \small\texttt{jpbald93@gmail.com}}
\date{27 September 2026\\ \small Version~2 adds the Layer-14 and Layer-15 checks (Section~\ref{sec:l1415})}

\emergencystretch=3em
\begin{document}
\maketitle

\begin{abstract}
We formalise, in Lean~4 with Mathlib, two classical elementary results about primitive roots
modulo primes of a special shape. The first is Chebyshev's criterion: if $p$ and $q=4p+1$ are
prime and $p\equiv 2\pmod 5$, then $10$ is a primitive root modulo~$q$. We also formalise the
converse, which follows from the same argument; we have not found it stated explicitly in the
literature. The second is the theorem of Ramesh and Makeshwari (2022) that the least positive
primitive root modulo a safe prime $p=2q+1$ is itself prime. Here the formal statement needs no separate existence hypothesis: we also prove that the least positive primitive root exists, using Mathlib's
cyclicity of $(\mathbb{Z}/p\mathbb{Z})^\times$. Both proofs reduce to the same criterion for primes $p=2^k\ell+1$ with $\ell$ an odd prime: a
nonzero residue is a primitive root exactly when it is a quadratic nonresidue and its $2^k$-th power
is not~$1$. We formalise this criterion for all $k\geq 1$, show that the two criteria above are equivalent, via
Euler's criterion, to its cases $k=2$ and $k=1$ (apart from the safe prime $5$), and derive three further classical families from it: $2$ is a primitive root
modulo every prime $4\ell+1$; $3$ is a primitive root modulo every prime $16\ell+1$; and $3$ is a
primitive root modulo every prime $8\ell+1$ except $41$. The development has seventeen main
theorems in about 700 lines. We
also describe a small checking script, stricter than a plain \texttt{lake build}, that rejects
several ways of faking a successful build. We test it against deliberately tampered versions of
the source. Version~2 adds two further checks. First, each of the seventeen theorems is restated with
Mathlib-only types and proved, with an explicit non-degenerate witness. Second, an AI writer
turned textbook statements of the results into Lean without seeing our code, and every one of its 26
statements was compared with the development. One of them, the general form of Wertheim's criterion as
transcribed, is false as written; a corrected form is proved.
\end{abstract}

\section{Introduction}

An integer $a$ is a \emph{primitive root} modulo a prime $p$ if its multiplicative order modulo
$p$ is exactly $p-1$. Artin's conjecture on primitive roots (see Moree's survey~\cite{Moree}) predicts
that every $a\neq -1$ that is not a perfect square is a primitive root for infinitely many
primes. Remarkably, this is not known for any single~$a$. Unconditional results come in two
flavours: statements for sets of $a$ (Heath-Brown~\cite{HeathBrown}), and \emph{explicit
families}, where a special shape of $p-1$ makes the multiplicative order easy to control.

This paper concerns two explicit families of the second kind:
\begin{enumerate}
\item \textbf{Chebyshev's criterion.} If $p$ and $q=4p+1$ are both prime and $p\equiv 2\pmod 5$,
then $10$ is a primitive root modulo~$q$. The statement is quoted in \cite[\S1]{Moree}. We
formalise it together with its converse.
\item \textbf{Ramesh--Makeshwari.} If $p=2q+1$ is a \emph{safe prime} (i.e.\ $q$ is also prime),
then the least positive primitive root modulo~$p$ is prime~\cite{RM}.
\end{enumerate}

Neither result is deep. We formalise them because each depends on the same short
piece of reasoning, and Mathlib now makes that reasoning inexpensive. There are three ingredients:
the prime-divisor test for the order of a group element, Euler's criterion, and quadratic reciprocity.
The formalisation also makes two small points explicit. One is the converse of Chebyshev's
criterion. The other is that we package the standard existence theorem for primitive roots with
the primality result, so the formal statement does not need a separate existence hypothesis.

\paragraph{Contributions.}
\begin{itemize}
\item A Lean~4 formalisation of both results, with seventeen named theorems (Table~\ref{tab:thms}), in the
namespace \lean{PrimitiveRootFamilies}.
\item A general criterion for primitive roots modulo primes $2^k\ell+1$
(Theorem~\ref{thm:gen}), stated once for all $k\geq1$, with the two families above as instances and
three further families as corollaries (Section~\ref{sec:general}).
\item A formal proof of the converse of Chebyshev's criterion
(Theorem~\ref{thm:cheb-iff}). It is elementary and very probably folklore. We do not claim it as new.
\item A formal form of the Ramesh--Makeshwari theorem with no separate existence hypothesis
(Theorem~\ref{thm:rm-exists}), with ``least'' expressed through Mathlib's \lean{IsLeast}, so
that the statement is about the least positive integer primitive root and not about a chosen
representative.
\item A proof-checking script (``gate'') and a set of tampering tests, described in
Section~\ref{sec:gate}.
\item (Version~2) A restatement of all seventeen theorems with Mathlib-only types, each with witnesses,
and a comparison with independently written classical statements (Section~\ref{sec:l1415}).
\end{itemize}

\paragraph{Artifact.} The development uses Lean~4~\cite{Lean4} and Mathlib~\cite{Mathlib}. The code is at \url{\repo}, commit \texttt{\commit}. It uses Lean
\texttt{v4.33.1} and Mathlib commit \texttt{0df444a3}, pinned in \texttt{lake-manifest.json}. Table~\ref{tab:thms} links each main declaration to that commit. The licence is Apache~2.0.
Version~2 of this paper adds the files \lean{PrimitiveRootFamilies/L14.lean} and
\lean{PrimitiveRootFamilies/L15Classical.lean} (Section~\ref{sec:l1415}) in a new commit,
\texttt{b4ae30a94f5bc5bcba7e0143c239565af261038a}. The statements and contents of \lean{Chebyshev.lean}, \lean{General.lean} and
\lean{SafePrime.lean} are unchanged; version~2 updates the root import file, the gate's
required-declaration list and the README. The links in
Table~\ref{tab:thms} still point to commit~\texttt{\commit}.
Supporting material (the Layer-14 and Layer-15 reports, the blind writer's statements and report, the transcript audit, and the reviews of this version with the fixes made in response) is in the folder \texttt{Primitive-root-families} of the repository \url{https://github.com/jpbald93/LEAN-verification}.

\begin{table}[t]
\centering\small
\begin{tabular}{@{}lll@{}}
\toprule
Lean name & Paper & File \\ \midrule
\lean{isPrimitiveRoot\_iff\_of\_eq\_four\_mul\_add\_one} & Lemma~\ref{lem:crit4} & \src{Chebyshev.lean}{L122} \\
\lean{legendreSym\_ten\_of\_eq\_four\_mul\_add\_one} & Lemma~\ref{lem:leg10} & \src{Chebyshev.lean}{L147} \\
\lean{isPrimitiveRoot\_ten\_iff\_mod\_five\_eq\_two} & Theorem~\ref{thm:cheb-iff} & \src{Chebyshev.lean}{L165} \\
\lean{isPrimitiveRoot\_ten\_of\_mod\_five\_eq\_two} & Corollary~\ref{cor:cheb} & \src{Chebyshev.lean}{L210} \\ \midrule
\lean{isPrimitiveRoot\_iff\_of\_eq\_two\_mul\_add\_one} & Lemma~\ref{lem:crit2} & \src{SafePrime.lean}{L54} \\
\lean{isPrimitiveRoot\_of\_legendreSym\_eq\_neg\_one} & Lemma~\ref{lem:nonres} & \src{SafePrime.lean}{L91} \\
\lean{prime\_of\_isLeast\_isPrimitiveRoot} & Theorem~\ref{thm:rm} & \src{SafePrime.lean}{L116} \\
\lean{exists\_isLeast\_isPrimitiveRoot} & Lemma~\ref{lem:exists} & \src{SafePrime.lean}{L181} \\
\lean{exists\_isLeast\_isPrimitiveRoot\_and\_prime} & Theorem~\ref{thm:rm-exists} & \src{SafePrime.lean}{L193} \\
\lean{isPrimitiveRoot\_two\_of\_mod\_four\_eq\_one} & Corollary~\ref{cor:two} & \src{SafePrime.lean}{L201} \\ \midrule
\lean{isPrimitiveRoot\_iff\_of\_eq\_two\_pow\_mul\_add\_one} & Theorem~\ref{thm:gen} & \src{General.lean}{L71} \\
\lean{isPrimitiveRoot\_of\_legendreSym\_eq\_neg\_one\_of\_not\_dvd} & Corollary~\ref{cor:gen-suff} & \src{General.lean}{L111} \\
\lean{isPrimitiveRoot\_two\_of\_eq\_four\_mul\_add\_one} & Corollary~\ref{cor:fam}(a) & \src{General.lean}{L120} \\
\lean{isPrimitiveRoot\_three\_of\_eq\_eight\_mul\_add\_one} & Corollary~\ref{cor:fam}(b) & \src{General.lean}{L167} \\
\lean{isPrimitiveRoot\_three\_of\_eq\_sixteen\_mul\_add\_one} & Corollary~\ref{cor:fam}(c) & \src{General.lean}{L204} \\
\lean{isPrimitiveRoot\_iff\_of\_eq\_two\_mul\_add\_one'} & Theorem~\ref{thm:gen}, $k=1$ & \src{General.lean}{L244} \\
\lean{isPrimitiveRoot\_iff\_of\_eq\_four\_mul\_add\_one'} & Theorem~\ref{thm:gen}, $k=2$ & \src{General.lean}{L252} \\
\bottomrule
\end{tabular}
\caption{The seventeen main theorems and where they are stated. Each depends only on Lean's standard
axioms \lean{propext}, \lean{Classical.choice} and \lean{Quot.sound}.}
\label{tab:thms}
\end{table}

\section{Conventions and the Mathlib interface}

Throughout, ``$a$ is a primitive root modulo $n$'' is Mathlib's
\lean{IsPrimitiveRoot (a : ZMod n) (n - 1)}. By definition, $a^{n-1}=1$, and $a^{k}=1$ implies
$(n-1)\mid k$. Equivalently (\lean{IsPrimitiveRoot.iff\_orderOf}), $a$ has multiplicative
order exactly $n-1$ in $\mathbb{Z}/n\mathbb{Z}$. For prime $n$ this is the standard notion.
Legendre symbols are Mathlib's \lean{legendreSym p a}, for $p$ prime and $a\in\mathbb{Z}$.

Everything rests on the following facts from Mathlib:
\begin{itemize}
\item \lean{orderOf\_eq\_of\_pow\_and\_pow\_div\_prime}:\\ if $n>0$, $x^n=1$ and $x^{n/r}\neq 1$ for every
prime $r\mid n$, then $\operatorname{ord}(x)=n$;
\item \lean{ZMod.pow\_card\_sub\_one\_eq\_one} (Fermat's little theorem);
\item \lean{legendreSym.eq\_pow} (Euler's criterion, for odd prime $p$): $\leg{a}{p}\equiv a^{(p-1)/2}\pmod p$;
\item \lean{legendreSym.at\_two} (the second supplementary law);
\item \lean{legendreSym.quadratic\_reciprocity\_one\_mod\_four} (reciprocity when one prime is $1\bmod 4$);
\item \lean{HasEnoughRootsOfUnity.exists\_primitiveRoot}, which provides a primitive $(p-1)$-th
root of unity in $\mathbb{Z}/p\mathbb{Z}$, i.e.\ cyclicity of the unit group.
\end{itemize}

\section{Chebyshev's base-10 family}

Let $p$ and $q=4p+1$ both be prime. Then $p\neq 2$, since $4\cdot 2+1=9$ is composite. So $p$ is
odd and $q\equiv 5\pmod 8$.

\begin{lemma}[Criterion for $q=4p+1$]\label{lem:crit4}
Let $p$ and $q=4p+1$ be prime and $a\not\equiv 0\pmod q$. Then $a$ is a primitive root modulo
$q$ if and only if $a^{2p}\not\equiv 1$ and $a^4\not\equiv 1\pmod q$.
\end{lemma}
\begin{proof}
The order of $a$ divides $q-1=4p$, whose prime divisors are $2$ and $p$. Apply the prime-divisor
test with $4p/2=2p$ and $4p/p=4$.
\end{proof}

\begin{lemma}\label{lem:leg10}
For primes $p$ and $q=4p+1$, $\leg{10}{q}=-\leg{q}{5}$.
\end{lemma}
\begin{proof}
Since $q\equiv 5\pmod 8$, $\leg{2}{q}=-1$. Since $5\equiv 1\pmod 4$, reciprocity gives
$\leg{5}{q}=\leg{q}{5}$. Multiply.
\end{proof}

\begin{theorem}[Chebyshev's criterion and its converse]\label{thm:cheb-iff}
Let $p$ and $q=4p+1$ be prime. Then $10$ is a primitive root modulo $q$ if and only if
$p\equiv 2\pmod 5$.
\end{theorem}
\begin{proof}
Note first that $10\not\equiv 0\pmod q$, because $q\geq 13$. Also $10^4\not\equiv 1\pmod q$:
otherwise $q\mid 9999=3^2\cdot 11\cdot 101$. The primes $3$ and $11$ are not of the form $4p+1$,
and $101=4\cdot 25+1$ with $25$ composite. By Euler's criterion, $10^{2p}\equiv\leg{10}{q}$.
Lemmas~\ref{lem:crit4} and~\ref{lem:leg10} then show that $10$ is a primitive root if and only if
$\leg{q}{5}=1$, i.e.\ $q\equiv 1$ or $4\pmod 5$. If $q\equiv 1\pmod 5$ then $5\mid p$, so
$p=5$ and $q=21$, which is not prime. And $q=4p+1\equiv 4\pmod 5$ exactly when $p\equiv 2\pmod 5$.
\end{proof}

\begin{corollary}[Chebyshev]\label{cor:cheb}
If $p$ and $q=4p+1$ are prime and $p\equiv 2\pmod 5$, then $10$ is a primitive root modulo $q$.
\end{corollary}

The formal statement of the corollary is
\begin{lstlisting}
theorem isPrimitiveRoot_ten_of_mod_five_eq_two {p q : ℕ} (hp : p.Prime) (hq : q.Prime)
    (hqp : q = 4 * p + 1) (hp5 : p % 5 = 2) : IsPrimitiveRoot (10 : ZMod q) (q - 1)
\end{lstlisting}
An earlier public version of the repository called this theorem \lean{ten\_isPrimitiveRoot}.
That name remains available as a deprecated alias.

\begin{remark}[On the converse]
The ``only if'' direction of Theorem~\ref{thm:cheb-iff} uses nothing beyond the argument for the
``if'' direction. We searched for an explicit statement and found none. The underlying criterion
is a standard textbook exercise: every quadratic nonresidue of a prime $4p+1$ is a primitive
root or has order~$4$ (cf.\ Burton~\cite[\S9.2, Problem~11(a)]{Burton}). So we regard the converse as
folklore and do not claim it as a new result.
\end{remark}

\paragraph{Formalisation notes.}
The main source of friction is that Mathlib's reciprocity and supplementary laws are stated for
natural-number primes cast to~$\mathbb{Z}$. Rewriting the goal into the form
\lean{legendreSym q (5 : $\mathbb{N}$)} took more steps than the mathematics suggests. Mathlib deliberately
provides no global instance \lean{Fact (Nat.Prime 5)}, so we declare one locally. The step
$q\nmid 9999$ is a four-way case split on the prime factors of $9999$. Each case is closed by
\lean{omega}, except for $101$, which needs the separate observation that $25$ is not prime. The
file ends with two \lean{example}s: $p=7$, $q=29$, where $10$ is a primitive root, and $p=3$,
$q=13$, where it is not.

\section{Safe primes}

Now let $p=2q+1$ with $p$ and $q$ both prime. The argument has the same structure, with the
exponent $4$ replaced by~$2$.

\begin{lemma}[Criterion for safe primes]\label{lem:crit2}
If $a\not\equiv 0\pmod p$, then $a$ is a primitive root modulo $p$ if and only if
$a^q\not\equiv 1$ and $a^2\not\equiv 1\pmod p$.
\end{lemma}

\begin{lemma}\label{lem:nonres}
Every quadratic nonresidue $x$ with $1<x<p-1$ is a primitive root modulo $p$.
\end{lemma}
\begin{proof}
By Euler's criterion $x^q\equiv\leg{x}{p}=-1\not\equiv 1$. And $x^2\equiv 1$ would force
$x\equiv\pm1\pmod p$, which the range $1<x<p-1$ excludes.
\end{proof}

\begin{theorem}[Ramesh--Makeshwari~\cite{RM}]\label{thm:rm}
Let $p=2q+1$ be a safe prime. If $g$ is the least element of
$\{a\in\mathbb{N}: a>0,\ a \text{ is a primitive root modulo } p\}$, then $g$ is prime.
\end{theorem}
\begin{proof}
First, $g<p$. Otherwise $g\bmod p$ would be a smaller positive primitive root. It is positive
because $p\nmid g$. By Lemma~\ref{lem:crit2} and Euler's criterion, $\leg{g}{p}=-1$. Hence
$g\neq 1$. Also $g\neq p-1$, since $(p-1)^2\equiv 1$. Suppose $g=mn$ with $1<m,n<g$. Then
$\leg{m}{p}\leg{n}{p}=-1$, so one of $m$ and $n$ is a nonresidue in $(1,p-1)$. By
Lemma~\ref{lem:nonres}, that factor is a primitive root smaller than $g$, a contradiction.
\end{proof}

Our proof is phrased with Legendre symbols. The published note~\cite{RM} argues with
multiplicative orders instead. The mathematical content is the same.

\paragraph{Existence.}
Theorem~\ref{thm:rm} is a statement about ``the least $g$''. It would be vacuously true if no
positive primitive root existed. We therefore also prove existence.

\begin{lemma}\label{lem:exists}
For every prime $p$, the set of positive integers that are primitive roots modulo $p$ has a least
element.
\end{lemma}
\begin{proof}
Mathlib provides a primitive $(p-1)$-th root of unity $\zeta\in\mathbb{Z}/p\mathbb{Z}$. Its
canonical representative $\zeta.\mathrm{val}\in\{0,\dots,p-1\}$ is nonzero and is a primitive
root. The least element then exists by well-ordering (\lean{Nat.find}).
\end{proof}

\begin{theorem}[With existence]\label{thm:rm-exists}
For every safe prime $p$, the least positive primitive root modulo $p$ exists and is prime.
\end{theorem}
\begin{lstlisting}
theorem exists_isLeast_isPrimitiveRoot_and_prime {p q : ℕ} [Fact p.Prime] (hq : q.Prime)
    (hpq : p = 2 * q + 1) :
    ∃ g, IsLeast {a : ℕ | 0 < a ∧ IsPrimitiveRoot (a : ZMod p) (p - 1)} g ∧ g.Prime
\end{lstlisting}

\begin{corollary}\label{cor:two}
If $p=2q+1$ is a safe prime with $q\equiv 1\pmod 4$, then $2$ is a primitive root modulo $p$.
\end{corollary}
\begin{proof}
Here $q\geq 5$ and $p\equiv 3\pmod 8$, so $\leg{2}{p}=-1$. Apply Lemma~\ref{lem:nonres}.
\end{proof}

\paragraph{Formalisation notes.}
The only step that is not purely routine in the formal proof of Theorem~\ref{thm:rm} is the bound
$g<p$. Informally, primitive roots are residues and nobody states it. In Lean the set is a set
of natural numbers, and the bound must be derived from minimality. We use
\lean{Nat.exists\_dvd\_of\_not\_prime2} to obtain a nontrivial factorisation $g=mn$, and
\lean{legendreSym.mul} for multiplicativity. Existence (Lemma~\ref{lem:exists}) takes ten lines,
because Mathlib already has the cyclicity of $(\mathbb{Z}/p\mathbb{Z})^\times$ in the form of
\lean{HasEnoughRootsOfUnity}.

\section{\texorpdfstring{The general criterion for primes $2^k\ell+1$}{The general criterion}}\label{sec:general}

Lemmas~\ref{lem:crit4} and~\ref{lem:crit2} are, via Euler's criterion and apart from the safe prime
$p=5$, equivalent to the cases $k=2$ and $k=1$ of a single criterion.
Throughout this section, $\ell$ is an odd prime, $k\geq 1$, and $p=2^k\ell+1$ is prime. Then
$p-1=2^k\ell$ has exactly the prime divisors $2$ and $\ell$.

\begin{theorem}\label{thm:gen}
Let $a\in\mathbb{Z}$ with $p\nmid a$. Then $a$ is a primitive root modulo $p$ if and only if
$\leg{a}{p}=-1$ and $a^{2^k}\not\equiv 1\pmod p$.
\end{theorem}
\begin{proof}
By the prime-divisor test, $a$ is a primitive root if and only if $a^{(p-1)/2}\not\equiv 1$ and
$a^{(p-1)/\ell}=a^{2^k}\not\equiv 1$. By Euler's criterion, $a^{(p-1)/2}\equiv\leg{a}{p}$, and
since $\leg{a}{p}=\pm1$ with $p>2$, the first condition is equivalent to $\leg{a}{p}=-1$.
\end{proof}

\begin{corollary}\label{cor:gen-suff}
If $\leg{a}{p}=-1$ and $p\nmid a^{2^k}-1$, then $a$ is a primitive root modulo $p$.
\end{corollary}
The formal version also assumes $a\not\equiv 0\pmod p$ explicitly. This is redundant, since it
follows from $\leg{a}{p}=-1$, but harmless.


The formal statement of Theorem~\ref{thm:gen} takes $a$ in $\mathbb{Z}$ rather than in
$\mathbb{Z}/p\mathbb{Z}$, so that the Legendre symbol can be stated directly:
\begin{lstlisting}
theorem isPrimitiveRoot_iff_of_eq_two_pow_mul_add_one {p l k : ℕ} [Fact p.Prime] (hl : l.Prime)
    (hl2 : l ≠ 2) (hk : 1 ≤ k) (hpl : p = 2 ^ k * l + 1) {a : ℤ} (ha : (a : ZMod p) ≠ 0) :
    IsPrimitiveRoot (a : ZMod p) (p - 1) ↔
      legendreSym p a = -1 ∧ (a : ZMod p) ^ (2 ^ k) ≠ 1
\end{lstlisting}
Two short bridge theorems (Table~\ref{tab:thms}) restate its cases $k=1$ and $k=2$ in this form.
The case $k=1$ needs the hypothesis $\ell\neq 2$, so it omits the safe prime $p=5$; the case
$k=2$ does not, since $4\cdot2+1=9$ is composite. The Layer-15 file of version~2 covers $p=5$ and the
textbook form with $a\not\equiv-1$ in place of $a^2\not\equiv1$ (Section~\ref{sec:l1415}). Lemmas~\ref{lem:crit4} and~\ref{lem:crit2}
were formalised first and are kept, with their original statements, for the proofs in the
preceding sections.

Theorem~\ref{thm:gen} is the case $p-1=2^k\ell$ of a sufficient criterion that Dickson attributes to
Wertheim~\cite[Vol.~I, p.~196]{Dickson}, together with its converse. As summarised there, the general
form for $p-1=2^k\ell^j$ is false when $j\geq2$ (Section~\ref{sec:l1415}).

The criterion turns a primitive-root question into two checks: a Legendre symbol, handled by
reciprocity, and finitely many exceptional primes, namely the prime divisors of $a^{2^k}-1$.

\begin{corollary}\label{cor:fam}
Let $\ell$ be prime.
\begin{enumerate}
\item[(a)] If $p=4\ell+1$ is prime, then $2$ is a primitive root modulo $p$.
\item[(b)] If $p=8\ell+1$ is prime and $\ell\neq 5$, then $3$ is a primitive root modulo $p$.
\item[(c)] If $p=16\ell+1$ is prime, then $3$ is a primitive root modulo $p$.
\end{enumerate}
\end{corollary}
\begin{proof}
(a) Here $\ell\neq 2$, since $9$ is composite, so $p\equiv 5\pmod 8$ and $\leg{2}{p}=-1$. If
$2^4\equiv1$ then $p\mid 15$, which is impossible since $p\geq 13$.

(b) The case $\ell=2$, $p=17$, is checked directly. Otherwise $\ell$ is odd and
$p\equiv 1\pmod 4$, so $\leg{3}{p}=\leg{p}{3}$. Now $p\equiv 2\ell+1\pmod 3$: $\ell\equiv1$
gives $3\mid p$, and $\ell=3$ gives $p=25$. Hence $p\equiv 2\pmod 3$ and $\leg{3}{p}=-1$.
If $3^8\equiv 1$, then $p\mid 6560=2^5\cdot5\cdot41$, so $p=41$ and $\ell=5$, which is excluded.

(c) Here $\ell\neq 2$, since $33$ is composite, and $p\equiv \ell+1\pmod 3$: $\ell\equiv 2$ gives
$3\mid p$, and $\ell=3$ gives $p=49$. So $\leg{3}{p}=-1$ as in (b). If $3^{16}\equiv1$ then
$p\mid 3^{16}-1=2^6\cdot5\cdot17\cdot41\cdot193$, and none of $5,17,41,193$ is of the form
$16\ell+1$ with $\ell$ prime.
\end{proof}

The exception in (b) is genuine: $3^8\equiv1\pmod{41}$, so $3$ is not a primitive root modulo
$41=8\cdot5+1$. The formal file proves this as an \lean{example}. In fact its order is $8$,
since $3^4\equiv-1\pmod{41}$. It illustrates the second of
the two checks: a nonresidue can fail to be a primitive root only at a prime divisor of
$a^{2^k}-1$. Corollary~\ref{cor:two} is a relative of part (a): it applies the general criterion with
$a=2$, $k=1$ and $\ell\equiv1\pmod 4$. Part (a) is Burton's \S9.2, Problem~11(b)~\cite{Burton}.
Parts (b) and (c) are elementary applications of the same method, and we make no novelty
claim for any of the three.

\section{Numerical cross-checks}

The repository contains three short Python scripts. They are independent of the Lean development and
serve as a sanity check that the formal statements say what we intend.
\begin{itemize}
\item \texttt{code/check\_family.py 200000}: there are $1916$ prime pairs $(p,4p+1)$ with
$p\leq 2\cdot 10^5$. In every case, $10$ is a primitive root modulo $4p+1$ exactly when
$p\equiv2\pmod5$. That happens for $628$ of the pairs. As a control, for $p=22$ (not prime) and
$q=89$, the order of $10$ is $44$.
\item \texttt{code/check\_safe\_primes.py 1000000}: there are $4324$ safe primes $p\leq 10^6$.
For every one, the least primitive root is prime. For the $2133$ with $q\equiv1\pmod4$, $2$ is a
primitive root. For the $72$ safe primes below $5000$, a brute-force computation of orders confirms
that all $73514$ nonresidues in $(1,p-1)$ are primitive roots. As a control, the least primitive
root of the prime $41$, which is not a safe prime, is the composite number~$6$.
\item \texttt{code/check\_general.py}: for the $16563$ primes $p=2^k\ell+1\leq 2\cdot10^6$ with
$\ell$ an odd prime, Theorem~\ref{thm:gen} agrees with SymPy's primitive-root test in all
$182191$ cases $(p,a)$ tested, for the bases $a\in\{2,3,5,6,7,10,11,12,-1,-2,-3\}$. SymPy's test
uses the same prime-divisor method, so as an independent control, orders are also computed by
brute force in the $1219$ cases with $p<3000$. In the same range, $2$ is a primitive root modulo all $4109$ primes $4\ell+1$, and
$3$ is a primitive root modulo all $1181$ primes $16\ell+1$. Among the $2158$ primes $8\ell+1$,
the only failure for $3$ is $p=41$.
\end{itemize}

For interest, Table~\ref{tab:dist} gives the distribution of the least primitive root over the
same $4324$ safe primes. The least primitive root is $2$ for the $2133$ primes with
$p\equiv 3\pmod 8$ and also for $p=5$. The value $3$ occurs only for $p=7$. This agrees with a
one-line computation, not formalised here. For a safe prime $p>7$, we have $p\equiv 3\pmod 4$
and $p\equiv 2\pmod 3$, so reciprocity gives $\leg{3}{p}=-\leg{p}{3}=+1$. Hence $3$ is never a
primitive root modulo such a~$p$.

\begin{table}[h]
\centering\small
\begin{tabular}{@{}l*{13}{r}@{}}
\toprule
$g$ & 2 & 3 & 5 & 7 & 11 & 13 & 17 & 19 & 23 & 29 & 31 & 37 & 41\\
count & 2134 & 1 & 1461 & 290 & 194 & 141 & 54 & 24 & 15 & 3 & 3 & 3 & 1\\
\bottomrule
\end{tabular}
\caption{Least primitive root $g$ of the $4324$ safe primes $p\leq 10^6$.}
\label{tab:dist}
\end{table}

\section{A tamper-resistant check}\label{sec:gate}

A \texttt{lake build} that succeeds does not by itself show that a statement has been proved.
The source may contain \lean{sorry}, a user-declared \lean{axiom}, or \lean{native\_decide}. It
may also contain code that runs at elaboration time and prints misleading output. For example,
\lean{\#eval} can print text that looks like a clean \lean{\#print axioms} report, and
\lean{macro\_rules} can redefine what \lean{\#print axioms} does. The development was prepared with
AI assistance (see Acknowledgements), so we treated these as real risks. The repository's
\texttt{gate.sh} therefore does the following:
\begin{enumerate}
\item It rejects any source file that contains \lean{sorry}, \lean{admit}, \lean{native\_decide},
the \lean{axiom} keyword, \lean{\#eval}, \lean{run\_cmd}, \lean{initialize}, \lean{IO},
\lean{set\_option}, or syntax-extension commands (\lean{macro}, \lean{elab}, \lean{syntax},
\lean{notation}, \lean{import Lean}, \ldots).
\item It builds the library.
\item It writes its own temporary file of \lean{\#print axioms} queries, one for each theorem in
Table~\ref{tab:thms}, instead of trusting any report produced inside the repository.
\item It flattens Lean's line-wrapped output, requires exactly one report per theorem and no
others, and accepts only the axioms \lean{propext}, \lean{Classical.choice} and \lean{Quot.sound}.
\end{enumerate}

We tested the gate by planting each of the following in a copy of the repository and confirming
that the gate fails: a \lean{sorry}; an \lean{axiom} whose keyword sits on a line of its own; a
fake clean report printed with \lean{\#eval}; a \lean{macro\_rules} redefinition of
\lean{\#print axioms}; an extra axiom hidden on a wrapped output line (with the source filter's \lean{axiom} check
switched off, so that the output parser is tested alone); and a false hypothesis
($p\bmod 5=3$ in place of $2$), which breaks the proof. The unmodified repository passes with
\texttt{PASS (17 theorems, standard axioms only)}. The script \texttt{tests/tamper.sh} in the artifact reproduces all six tests: each one plants its
defect in a scratch copy of the repository and checks that the gate fails at the expected stage.

\subsection{Layers 14 and 15}\label{sec:l1415}
Version~2 of the artifact adds two more checks, called Layers~14 and~15. Both were carried out by AI
agents. The library files \lean{Chebyshev.lean}, \lean{General.lean} and \lean{SafePrime.lean} are
unchanged: their hashes were compared before and after. Version~2 updates the root import file
and the gate's required-declaration list.

\paragraph{Layer 14: Mathlib-only statements and witnesses.}
\begin{sloppypar}The file \lean{PrimitiveRootFamilies/L14.lean} (376 lines) restates all seventeen theorems of
Table~\ref{tab:thms}. No theorem is omitted. The published statements already mention only Mathlib
constants, so each restatement goes further and uses textbook phrasings. Primality is an explicit
hypothesis \lean{Nat.Prime p}, not an instance. ``Primitive root'' is \lean{orderOf a = p - 1}, via
\lean{IsPrimitiveRoot.iff\_orderOf}. The Legendre symbol is \lean{jacobiSym a p}, via
\lean{jacobiSym.legendreSym.to\_jacobiSym}. Each restatement \lean{l14\_$\ast$} is proved from the
library theorem and these bridges, with the same hypotheses. A mechanical check of the elaborated
types found no \lean{PrimitiveRootFamilies} constant in any of the 36 new declarations.\end{sloppypar}

Each restatement also has a witness theorem \lean{l14w\_$\ast$}, which proves every hypothesis on an
explicit instance together with a non-degenerate conclusion. For each equivalence, the witness
exhibits both truth values. Two of them show that the condition $a^{2^k}\not\equiv1$ in
Theorem~\ref{thm:gen} is needed: $a=3$ for $p=41$, and $a=5$ for $p=13$. In each case $a$ is a
nonresidue but not a primitive root. For Theorem~\ref{thm:rm} the hypothesis \lean{IsLeast} is
realised explicitly: the least primitive root of $23$ is~$5$. For Corollary~\ref{cor:fam}(b) the
witness uses the generic case $p=89$, not just $p=17$, and also proves that $3$ is not a primitive
root modulo~$41$. No Mathlib-only statement turned out to be false or to need an extra hypothesis.
With Layer~14 the gate passes with \texttt{PASS (53 theorems, standard axioms only)}, and the
tampering tests still reject every planted defect. One test (\texttt{wrap}) replaces a library
theorem by \lean{True}, so a downstream file that used that theorem stopped building and the test
failed at the wrong stage. The restatement of Corollary~\ref{cor:cheb} is therefore derived from
Theorem~\ref{thm:cheb-iff}.

\paragraph{Layer 15: comparison with classical statements.}
A separate agent collected the classical statements from the literature, with citations: Burton,
Hardy--Wright, Dickson's \emph{History}~\cite{Dickson}, the Ramesh--Makeshwari note~\cite{RM}, and
encyclopaedic sources. They fall into six groups: Euler's criterion; the primitive-root test; the
Fermat-prime, safe-prime and Wertheim characterisations; Chebyshev's theorems on $2$ and on $10$;
and Ramesh--Makeshwari, with the lemma that the least nonresidue is prime. An AI writer received only
these statements and saw no code and no paper. Its transcript was audited, and its
report records that it opened no other file. It
wrote 26 Lean statements with Mathlib notions, including two existence obligations, all
with \lean{sorry}. A comparer preserved their hypotheses and conclusions, repackaging the theorem declarations as universally quantified proposition definitions, and proved each one from the development and
Mathlib, or refuted it. No hypothesis was added and the library was not changed. All 28 resulting
theorems depend only on the three standard axioms. The comparison was first checked as a scratch
file; it was then integrated as \lean{PrimitiveRootFamilies/L15Classical.lean} (704 lines), in the
namespace \lean{PrimitiveRootFamilies.Classical}, with the definitions and proposition statements
byte-identical to the checked comparison file and every proof kept. With both layers the gate passes with
\texttt{PASS (83 theorems, standard axioms only)}, and all six tampering tests are still rejected.
The outcome for the 26 statements:
\begin{itemize}
\item \emph{Covered} (8): the statements on $2$ modulo $4q+1$ (the library version does not need
$q\neq2$), the three forms of Chebyshev's base-10 criterion, its examples $29,149,269$, the
note that $p$ is odd and $q\equiv5\pmod8$, and Ramesh--Makeshwari with its existence obligation.
\item \emph{Narrower} (1): the safe-prime criterion ``primitive root $\iff\leg{a}{p}=-1$ and
$a\not\equiv-1$''. The library has it only for $\ell\neq2$ and with $a^2\not\equiv1$, so the case
$p=5$ and the translation to $a\not\equiv-1$ were missing. Both are supplied by a new proof, so the
statement is now proved in full.
\item \emph{Out of scope} (16): background results that the paper uses from Mathlib or does not claim.
These are Euler's criterion (3 statements), the primitive-root test and its corollaries (4), Fermat
primes, two remarks on $-1$ for safe primes, the decimal-period form of Chebyshev's criterion, the
examples of least primitive roots, the lemma that the least nonresidue is prime and exists (2), and
``least primitive root $=$ least nonresidue'' for safe primes with its bound (2). All 16 were proved.
\item \emph{Writer statement false as written} (1): Wertheim's criterion for $p=2^k q^l+1$ with
$k,l\geq1$. It is false for $l\geq2$: for $p=19=2\cdot3^2+1$ and $a=8$ we have $\leg{8}{19}=-1$ and
$8^2\not\equiv1$, but $8$ has order~$6$. The writer transcribed the source statement faithfully, so
the error lies in the source. Theorem~\ref{thm:gen} is its case $l=1$, which is correct and is
derived from the library. A corrected criterion, with $a^{2^kq^{l-1}}\not\equiv1$, is proved.
\end{itemize}
These proofs were absent from the version-1 library and are now included in the Layer-15 file:
\begin{itemize}
\item \lean{cov\_Prime3\_b}: the safe-prime criterion including $p=5$, with the condition $a\not\equiv-1$.
\item \lean{not\_classical\_Prime3\_c}: the counterexample $p=19$, $a=8$.
\item \lean{new\_Prime3\_c\_corrected}: Wertheim's criterion for general $l$, corrected.
\item \lean{cov\_Prime3\_a}: for a Fermat prime, the primitive roots are exactly the nonresidues.
\item \lean{cov\_Prime5\_reptend}: for primes $p$ and $q=4p+1$, $10$ is a primitive root modulo $q$
iff the decimal period of $1/q$ is $q-1$, using \lean{decimalPeriod\_eq\_orderOf}, which proves the
period/order equality for prime $q$ with $q\nmid10$.
\item \lean{cov\_Prime6\_a}, \lean{cov\_Prime6\_a\_obligation}: the least positive nonresidue exists and is prime.
\item \lean{cov\_Prime6\_b\_aux}, \lean{cov\_Prime6\_b}: for a safe prime, the least nonresidue is
below $p-1$ and equals the least primitive root.
\item thin wrappers around Mathlib for Euler's criterion and the primitive-root test.
\end{itemize}

\paragraph{Limitations.}
The token filter is a heuristic over the project's own sources, and it is deliberately
over-strict. The gate trusts the pinned toolchain, Mathlib and lake configuration. At commit~\texttt{\commit} it
audits only the seventeen named theorems; with Layer~14 it audits 53 declarations, and with
Layers~14 and~15 together 83. It is not a sandbox. It also does not compare theorem statements against a trusted specification. Layer~15 compares
them with independently written classical statements, but that comparison was done by AI agents. That
the formal statements mean what the paper says still requires a human reading, which
Sections~3--5 are intended to support. Environment replay with \texttt{leanchecker} (formerly
\texttt{lean4checker}) rechecks compiled declarations using Lean's own kernel, and independently
implemented external checkers provide a different trust boundary. Our gate is a lightweight
complement to these tools, not a substitute.

\section{Related work}

Mathlib contains the general theory used here: primitive roots, the cyclicity of
$(\mathbb{Z}/p\mathbb{Z})^\times$, Euler's criterion and quadratic reciprocity. We found neither family in the pinned Mathlib version. We found no prior formalisation of either
result in Lean, Isabelle, Coq/Rocq or Mizar, but our search was not exhaustive. The broader
context, Artin's conjecture and its partial results, is surveyed by Moree~\cite{Moree}. Other
formalisations of elementary number theory include the Isabelle AFP and Mathlib's
\texttt{NumberTheory} library.

\section{Conclusion}

Chebyshev's criterion, with its converse, and the Ramesh--Makeshwari theorem both follow from one
prime-divisor criterion together with Euler's criterion. With current Mathlib, each can be
formalised in about 220 lines, and the criterion itself holds uniformly for primes $2^k\ell+1$ (Section~\ref{sec:general}). Formalising them made two points explicit: the converse of Chebyshev's criterion, and the
existence of the least root. Natural extensions are primes $p$ with $p-1=2^k\ell^j$ or with more prime factors, where the
exponents change and each additional distinct prime factor adds a condition. The sufficient direction
for $2^k\ell^j$, with the condition $a^{2^k\ell^{j-1}}\not\equiv1$, is proved in the Layer-15 file;
the converse and further prime factors remain. Other extensions are criteria that also use cubic reciprocity (Fueter's
results, surveyed in~\cite[\S9.7, item~33]{Moree}), and a formal proof of the observation in Table~\ref{tab:dist} about the least primitive root.

\section*{Acknowledgements and use of AI}
The formal proofs, the checking script and a draft of this paper were prepared with the help of
large-language-model AI assistants. Two different models independently reviewed the
formalisation. The author checked each formal statement in Table~\ref{tab:thms} against its
informal counterpart by hand. The Layer-14 and Layer-15 checks of version~2
(Section~\ref{sec:l1415}) were carried out entirely by AI agents. Every proof is checked by the Lean kernel, and the author takes responsibility
for all content.

\begin{thebibliography}{9}
\bibitem{Burton} D.~M. Burton, \emph{Elementary Number Theory}, 7th ed., McGraw--Hill, 2011.
\bibitem{Dickson} L.~E. Dickson, \emph{History of the Theory of Numbers}, Vol.~I, Carnegie
Institution, Washington, 1919.
\bibitem{HeathBrown} D.~R. Heath-Brown, Artin's conjecture for primitive roots,
\emph{Quart.\ J. Math.\ Oxford} (2) \textbf{37} (1986), 27--38.
\bibitem{Mathlib} The mathlib Community, The Lean mathematical library, in \emph{Proc.\ 9th ACM
SIGPLAN Int.\ Conf.\ on Certified Programs and Proofs (CPP 2020)}, 367--381.
\bibitem{Lean4} L.~de Moura and S.~Ullrich, The Lean~4 theorem prover and programming language,
in \emph{Automated Deduction -- CADE 28}, LNCS \textbf{12699}, Springer, 2021, 625--635.
\bibitem{Moree} P.~Moree, Artin's primitive root conjecture -- a survey, \emph{Integers}
\textbf{12} (2012), no.~6, 1305--1416. doi:10.1515/integers-2012-0043.
\bibitem{RM} V.~P. Ramesh and M.~Makeshwari, Least primitive root of any safe prime is prime,
\emph{Amer.\ Math.\ Monthly} \textbf{129} (2022), no.~10, 971. doi:10.1080/00029890.2022.2115816.
\end{thebibliography}

\end{document}
